% =====================================================================
%  GMN2 §2  综述（Review）
%  原文：wkb-paper.tex 1032–1373 行；PDF（印刷）p.14–18
%  编号公式 (2.1)–(2.11) 共 11 条；脚注 4 条
%  块内公式用 \eqref；跨块引用暂写原文节号并加 % TODO
% =====================================================================

\section{综述}
\label{sec:review}

我们先把 \cite{Gaiotto:2008cd} 的设定、记号与主要提议扼要回顾一遍；更详细的讨论请参看该文。

\subsection{设定}

我们考虑一个 $d=4$、$\CN=2$ 超对称规范理论，记其库仑分支为 $\CB$。在每一点 $u \in \CB$ 处，规范群破缺到极大环面 $U(1)^r$。有一个电荷格 $\hat\Gamma_u$，其上带一个反对称的整值配对 $\inprod{,}$。此配对的根是味荷子格 $(\Gamma_{\rm flavor})_u$。对它取商，便得到规范荷的商格 $\Gamma_u$。$\Gamma_u$ 的秩为 $2r$，并带一个辛配对。

格 $\hat\Gamma_u$ 是某个局部系统 $\hat\Gamma$ 的纤维；围绕 $\CB$ 中那些复余维一的奇点轨迹——在那些地方，某些 BPS 粒子变成无质量——$\hat\Gamma$ 有非平凡的单值。存在一个「中心荷」同态
\begin{equation}
Z: \hat\Gamma_u \to \IC \label{ch2:eq:central-charge}
\end{equation}
它随 $u$ 全纯地变化。特别地，给定 $\hat\Gamma$ 的任一局部截面 $\gamma$，就有相应的、局部定义的全纯函数 $Z_\gamma(u)$，即电荷为 $\gamma$ 的粒子的中心荷。

我们把该理论表述在 $\IR^3 \times S^1$ 上，其中 $S^1$ 的半径为 $R$。在能量 $\ll 1/R$ 时，这个理论看起来是有效三维的。它的模空间局部上是两个\hk 流形之积。一个因子是 $d=4$ 理论的 Higgs 分支，本文不予考虑；另一个因子是 $d=3$ 的库仑分支 $\CM$。在 $\CM$ 的一般非奇异点处，Higgs 分支实际上并不出现。$\CM$ 是一个纤维化
\begin{equation}
\pi: \CM \to \CB \label{ch2:eq:fibration}
\end{equation}
其一般纤维是 $2r$-环面。环面纤维之所以出现，是因为 $d=4$ 中的规范场在 $d=3$ 中给出标量，即规范场（电与磁）绕 $S^1$ 的和乐（holonomy）。对每个 $\gamma \in \Gamma_u$，我们有相应的圆值和乐 $\theta_\gamma$，满足 $\theta_{\gamma+\gamma'} = \theta_\gamma + \theta_{\gamma'}$。

由于超对称，$\CM$ 上的度量 $g$ 是超凯勒的。对 $g$ 的一个初步近似 $g^{\sf}$ 可由朴素的量纲约化得到。要精确地确定 $g$，还必须计入瞬子（以及多瞬子）效应，它们来自 $d=4$ 理论中绕 $S^1$ 缠绕的 BPS 粒子。这些效应的权重是第二螺旋度超迹 $\Omega(\gamma;u)$，后者计数电荷为 $\gamma$ 的粒子。

\subsection{达布坐标}

在 \cite{Gaiotto:2008cd} 中我们给出了 $g$ 的一个精确描述。主要想法是：要描述 $g$，只需描述 $\CM$（看作全纯辛流形）上的全纯达布坐标。

这一构造在底空间 $\CB$ 上是局部的。固定一个开集 $U \subset \CB$，使 $\hat\Gamma$ 在其上可平凡化；再固定反对称配对模 $2$ 的一个二次加细 $\sigma: \hat\Gamma \to \{ \pm 1 \}$。我们的全纯达布坐标由 $\hat\Gamma$ 在 $U$ 上的截面 $\gamma$ 来标号。它们是 $\pi^{-1}(U) \times \IC^\times$ 上的函数 $\CX_\gamma$，满足

\begin{enumerate}[label=\alph*)]

\item $\CX_{\gamma+\gamma'} = \CX_\gamma \CX_{\gamma'}$.

\item 对任意固定的 $\zeta \in \IC^\times$，$\CX_\gamma(\cdot; \zeta)$ 取值于 $\IC^\times$，并且关于 $\CM$ 上的复结构 $J^{\pz}$ 全纯。（回忆一下：$\CM$ 是\hk 的，因而有一族 $\IC\IP^1$ 的复结构。）

\item 诸 $\CX_\gamma$ 的全纯泊松括号由下式给出\footnote{为后文方便，我们相对 \cite{Gaiotto:2008cd} 把泊松括号重新标度了一个因子 $4\pi^2 R$。}：
\begin{equation} \label{ch2:eq:poisson}
\left\{ \CX_\gamma, \CX_{\gamma'} \right\} = \inprod{\gamma, \gamma'} \CX_\gamma \CX_{\gamma'}.
\end{equation}

\item 对任意固定的 $(u, \theta) \in \CM$，$\CX_\gamma(u, \theta; \zeta)$ 关于 $\zeta$ 全纯。这里 $\theta$ 是 $\CM \to \CB$ 在 $u$ 上方那条纤维上的角坐标。

\item $\CX_{\gamma}(\cdot; \zeta) = \overline{\CX_{-\gamma}(\cdot; -1/\bar\zeta)}$.

\end{enumerate}

此外，在 \cite{Gaiotto:2008cd} 中，考虑还满足如下进一步渐近条件的坐标系显得格外有意思，即

\begin{enumerate}[label=\alph*)]
\setcounter{enumi}{5}

\item $\lim_{\zeta \to 0} \CX_\gamma(u, \theta; \zeta) \exp \left[ - \zeta^{-1} \pi R Z_\gamma(u) \right]$ 存在。

\end{enumerate}

然而，这里有一个斯托克斯现象在起作用：对于我们所关心的 $\CM$，结果是不可能构造出满足条件 a)–f) 的达布坐标。正确的做法是：只要求一切所需性质在某个以射线 $e^{i \vartheta}\IR_+$ 为中心的半平面内成立\footnote{我们要强调：在某些情形下，$\CX_\gamma^\vartheta$ 实际上依赖于 $\vartheta \in \IR$，而不只是 $\vartheta \in \IR/2\pi\IZ$。下文第 8 节% TODO 合稿改 \ref{sec:irregular}
、第 9 节% TODO 合稿改 \ref{section:Superconformal}
与第 11 节% TODO 合稿改 \ref{sec:UniversalStokesMatrix}
会遇到这种情形。}，
\begin{equation} \label{ch2:eq:halfplane}
\IH_\vartheta := \left\{ \zeta: \vartheta-\frac{\pi}{2} < \arg \zeta < \vartheta+\frac{\pi}{2} \right\},
\end{equation}
并且只对单个 $u_0 \in U \subset \CB$ 成立。于是我们要求的是一\emph{族}坐标系 $\CX_\gamma^{\vartheta, u_0}$，每一个都定义在 $\pi^{-1}(U) \times \IH_\vartheta$ 上。每一个都应满足 a)–c)，以及

\begin{enumerate}[label=\alph*')]
\setcounter{enumi}{3}

\item $\CX^{\vartheta, u_0}_\gamma(u, \theta; \zeta)$ 关于 $\zeta$ 全纯，只要 $\zeta \in \IH_\vartheta$。

\item $\CX^{\vartheta, u_0}_{\gamma}(\cdot; \zeta) = \overline{\CX^{\vartheta+\pi, u_0}_{-\gamma}(\cdot; -1/\bar\zeta)}$.

\item $\lim_{\zeta \to 0} \CX_\gamma^{\vartheta, u_0}(u_0, \theta; \zeta) \exp \left[ - \zeta^{-1} \pi R Z_\gamma(u_0) \right]$ 存在，其中 $\zeta$ 限制在 $\IH_\vartheta$ 内。

\end{enumerate}

本文将对充分大的 $R$ 给出满足这些条件的函数 $\CX_\gamma^{\vartheta, u_0}$ 的一个构造。这些 $\CX_\gamma^{\vartheta, u_0}$ 确实依赖于 $(\vartheta, u_0)$：在 $(\vartheta, u_0)$ 的空间中存在一个实余维一的子集，$\CX_\gamma^{\vartheta, u_0}$ 在其上发生跳变。由 b)、c) 可知，这些跳变都是全纯泊松态射。

这些坐标 $\CX_\gamma^{\vartheta, u_0}$ 的跳变是整个故事中最重要的部分。特别地，跳变由 $d=4$ 场论的 BPS 简并度所决定，并且反过来决定后者。若把 $u_0$ 视为固定，则跳变发生在特定的 $\vartheta$ 值处，即那些等于以 $u_0$ 标号的真空里 BPS 态中心荷相位的 $\vartheta$。此外，跳变的具体形式由 BPS 态的规范荷决定。我们更精确地表述如下：

\begin{enumerate}[label=\alph*)]
\setcounter{enumi}{6}

\item $\CX_\gamma^{\vartheta, u_0}$ 作为 $(\vartheta, u_0)$ 的函数是分段常值的，其不连续点出现在这样的数对 $(\vartheta, u_0)$ 处：存在某个 $\gamma_{\rm BPS}$，使得 $\arg -Z_{\gamma_{\rm BPS}}(u_0) = \vartheta$ 且 $\Omega(\gamma_{\rm BPS}; u_0) \neq 0$。

\item 固定 $\vartheta_0 \in \IR/2\pi\IZ$、$u_0 \in \CB$，并定义
\begin{equation} \label{ch2:eq:Sdef}
\mathbf{S}_{\vartheta_0, u_0} := \prod_{\gamma_{\rm BPS}: \ \arg -Z_{\gamma_{\rm BPS}}(u_0) = \vartheta_0} \CK_{\gamma_{\rm BPS}}^{\Omega(\gamma_{\rm BPS}; u_0)},
\end{equation}
其中 $\CK_{\gamma_{\rm BPS}}$ 是诸 $\CX_\gamma$ 的一个全纯泊松变换，由
\begin{equation} \label{ch2:eq:KStransf}
\CK_{\gamma_{\rm BPS}}: \CX_{\gamma} \mapsto \CX_{\gamma} \left( 1-\sigma(\gamma_{\rm BPS}) \CX_{\gamma_{\rm BPS}} \right)^{\inprod{\gamma, \gamma_{\rm BPS}}}
\end{equation}
给出。于是\footnote{我们这里的约定与 \cite{Gaiotto:2008cd} 相差一个符号。}
\begin{equation} \label{ch2:eq:jump}
\left( \lim_{\vartheta \to \vartheta_0^+} \CX^{\vartheta, u_0}_\gamma \right) = \mathbf{S}_{\vartheta_0, u_0} \left( \lim_{\vartheta \to \vartheta_0^-} \CX^{\vartheta, u_0}_\gamma \right).
\end{equation}
（这里作了一个一般性假设：满足 $\arg -Z_{\gamma_{\rm BPS}}(u_0) = \vartheta_0$ 的那些电荷都彼此成比例，因此无需对 $\mathbf{S}_{\vartheta_0, u_0}$ 中的乘积排序。）

\end{enumerate}

最后一个重要性质是

\begin{enumerate}[label=\alph*)]
\setcounter{enumi}{8}

\item 当 $\zeta \in \IH_\vartheta$ 时，有如下大 $R$ 渐近
\begin{equation} \label{ch2:eq:largeR}
\CX_{\gamma} \sim \CX^{\sf}_{\gamma},
\end{equation}
而且事实上修正项是指数式小的，即在那些与 $\CB$ 的奇点保持正距离的区域内，$\CX_\gamma = \CX^{\sf}_\gamma (1 + \CO(e^{-{\rm const} \cdot R}))$。

\end{enumerate}

在本文余下的部分，我们有时会放松记号：在不刻意强调 $u_0$ 依赖时，把 $\CX_\gamma^{\vartheta, u_0}$ 简写作 $\CX_\gamma^\vartheta$，甚至简写作 $\CX_\gamma$。（这个 $\CX_\gamma$ \emph{不}可与 \cite{Gaiotto:2008cd} 中的 $\CX_\gamma$ 混淆：那些函数在本文中记作 $\CX^{\RH}_\gamma$。见下文第 \ref{subsec:RH-problem} 节。）

\subsection{穿墙}
\label{subsec:wallcrossing}

正如我们刚刚回顾的，至少对充分大的 $R$，模空间 $\CM$ 上带有一族局部坐标系 $\CX^{\vartheta, u_0}_\gamma$，满足条件 a)–c)、d$'$)–f$'$)、g)–i)。

这些坐标的仅仅是存在，就已经给出相当强的推论。考虑一点 $u \in \CB$ 与两个不同的相位 $\vartheta_\pm$，满足 $\vartheta_+ - \vartheta_- < \pi$。坐标系 $\CX^{\vartheta_\pm, u}_\gamma$ 一般并不相等；要看出它们如何相关联，须对介于 $\vartheta_-$ 与 $\vartheta_+$ 之间的每条斯托克斯线应用 \eqref{ch2:eq:jump}。由此得到关系
\begin{equation} \label{ch2:eq:relate}
\CX_{\gamma}^{\vartheta_+, u} = \mathbf{S}(\vartheta_-, \vartheta_+; u) \CX_{\gamma}^{\vartheta_-, u},
\end{equation}
其中
\begin{equation} \label{ch2:eq:product}
\mathbf{S}(\vartheta_-, \vartheta_+; u) = \prod^{\ccwarrow}_{\gamma_{\rm BPS}: \ \vartheta_- < \arg -Z_{\gamma_{\rm BPS}}(u) < \vartheta_+} \CK_{\gamma_{\rm BPS}}^{\Omega(\gamma_{\rm BPS}; u)},
\end{equation}
而乘积按 $\arg -Z_{\gamma_{\rm BPS}}(u)$ 递增的次序取。

一个关键的代数事实 \cite{ks1} 是：形如 \eqref{ch2:eq:product} 的乘积分解是唯一的——\eqref{ch2:eq:product} 实际上\emph{确定}了 $\arg -Z_{\gamma_{\rm BPS}}(u)$ 介于 $\vartheta_-$ 与 $\vartheta_+$ 之间的那些 $\Omega(\gamma_{\rm BPS}; u)$（例外是 $\gamma_{\rm BPS} \in \Gamma_{\rm flavor}$，它们有 $\CK_{\gamma_{\rm BPS}} = \mathrm{id}$，因而在 $\mathbf{S}(\vartheta_-, \vartheta_+; u)$ 中不可见）。于是可以把 $\mathbf{S}(\vartheta_-, \vartheta_+; u)$ 看作诸 $\Omega(\gamma_{\rm BPS}; u)$ 的某种「生成函数」。

现在设我们把 $u$ 连续形变到 $u'$。只要在这个过程中没有 $\arg Z_{\gamma_{\rm BPS}}(u)$ 越过 $\vartheta_+$ 或 $\vartheta_-$，由 g) 即得 $\CX_\gamma^{\vartheta_\pm, u} = \CX_\gamma^{\vartheta_\pm, u'}$，从而
\begin{equation} \label{ch2:eq:wcf}
\mathbf{S}(\vartheta_-, \vartheta_+; u) = \mathbf{S}(\vartheta_-, \vartheta_+; u').
\end{equation}
公式 \eqref{ch2:eq:wcf} 在给定 $\Omega(\gamma_{\rm BPS}; u)$ 之后确定 $\Omega(\gamma_{\rm BPS}; u')$，因而给出了穿墙问题的一个完整解答。在我们的语境中，它只是函数 $\CX^{\vartheta, u}_\gamma$ 存在性的直接推论。

由于乘积 \eqref{ch2:eq:product} 一般是无穷的，我们或许应当说明它该如何理解。首先要找到 $\hat\Gamma$ 的某一组基 $\{\gamma_i\}$，使得对 $\mathbf{S}(\vartheta_-, \vartheta_+; u)$ 有贡献的所有 $\gamma_{\rm BPS}$ 都是诸 $\gamma_i$ 的\emph{非负}线性组合。这组基依赖于 $(\vartheta_-, \vartheta_+, u)$。\footnote{其实，这样的基的存在性并不显然。在 \cite{ks1} 的数学工作中附加了一个技术性条件来保证它，人们希望该条件在所有物理例子中都成立。本文第 7.8 节% TODO 合稿改 \ref{sec:bps-comments}
将证明，在我们所考虑的例子中这样的基确实存在。}每个 $\CK_\gamma$ 定义了代数 $F := \IC[[x_{\gamma_1}, \dots, x_{\gamma_i}]]$ 的一个泊松自同构。（无穷泰勒级数出现，是因为我们展开分母：$1/(1-x) = 1 + x + \cdots$。）那么它们的乘积如何理解？为此，我们先定义一族有限维幂幺群 $G_N(\vartheta_-, \vartheta_+; u)$（$N \ge 0$）。先用基 $\gamma_i$ 把 $\gamma = \sum n_i \gamma_i$ 的次数定义为 $\vert \gamma \vert := \sum_i n_i$，并粗略地说「把 $F$ 截断到次数不超过 $N$ 的 Fourier 模」。更精确地说，$F$ 有一列由次数大于 $N$ 的单式生成的理想 $I_N \subset F$ 给出的滤过。由于 $\CK_\gamma$ 把 $I_N$ 映入自身，满足 $\vartheta_- < -\arg Z_\gamma(u) < \vartheta_+$ 的那些变换生成 $F_N := F / I_N$ 的一个泊松自同构群 $G_N(\vartheta_-, \vartheta_+; u)$。此外，仍因诸 $\CK_\gamma$ 保持该滤过，存在投影映射 $G_N \to G_{N-1}$；用这些映射可把群 $G(\vartheta_-, \vartheta_+; u)$ 定义为诸 $G_N(\vartheta_-, \vartheta_+; u)$ 所成系统的逆向极限。投影到任一 $G_N(\vartheta_-, \vartheta_+; u)$ 之后，乘积 \eqref{ch2:eq:product} 只涉及有限多个非平凡因子，因而良定义；并且它与投影 $G_N \to G_{N-1}$ 相容。这就足以在 $G(\vartheta_-, \vartheta_+; u)$ 中定义它。

我们再来稍多谈谈上面所称的乘积分解的唯一性。有 $\CK_\gamma = e^{f_\gamma}$，其中 $f_\gamma := \sum_{n \ge 1} \frac{\sigma(n\gamma)}{n^2} \{ \CX_{n \gamma}, \cdot \}$。满足 $\vartheta_- < -\arg Z_\gamma(u) < \vartheta_+$ 的诸 $f_\gamma$ 张成 $G(\vartheta_-, \vartheta_+; u)$ 的李代数。对幂幺群 $G_N$，指数映射是双射。由此可知，一个群元素 $g \in G(\vartheta_-, \vartheta_+; u)$ 分解成形如 \eqref{ch2:eq:product} 的乘积确实唯一。从更实用的角度看，有一个从 $\mathbf{S}(\vartheta_-, \vartheta_+; u)$ 中提取诸 $\Omega(\gamma_{\rm BPS}; u)$ 的算法，对任何合理小的 $\gamma_{\rm BPS}$ 都容易在计算机上实现。它归结为归纳地考虑 $\mathbf{S}(\vartheta_-, \vartheta_+; u)$ 到各层子群 $G_N(\vartheta_-, \vartheta_+; u)$ 的投影：在第 $N$ 步，可以定出满足 $\vert \gamma_{\rm BPS} \vert \le N$ 的 $\Omega(\gamma_{\rm BPS}; u)$。

公式 \eqref{ch2:eq:wcf} 最早在 \cite{ks1} 中以非常一般的框架给出。它推广了原始与半原始穿墙公式，后者最先由 \cite{Denef:2007vg} 利用 Denef 的 $\CN=2$ 超引力的多中心 BPS 黑洞解 \cite{Denef:2000nb} 导出。\cite{Denef:2007vg} 的论证对超引力及其场论极限都适用。遗憾的是，\cite{Gaiotto:2008cd} 与本文的构造都只限于场论。于是仍留下一个重要公开问题：对紧 Calabi–Yau 上的 II 型弦理论给出 \eqref{ch2:eq:wcf} 的物理推导。

\subsection{黎曼–希尔伯特问题}
\label{subsec:RH-problem}

在 \cite{Gaiotto:2008cd} 中，我们并未显式引入函数 $\CX^\vartheta_\gamma$；我们表述的是一个黎曼–希尔伯特问题，其解将给出\emph{单个}达布坐标系 $\CX^{\RH}_\gamma(\zeta)$，它对一切 $\zeta$ 满足 a)–f)，\emph{例外}仅在于：$\CX^{\RH}_\gamma(\zeta)$ 关于 $\zeta$ 不是全纯的，而只是分片全纯的；它沿每条射线 $\zeta \in Z_{\gamma_{\rm BPS}}\IR_-$ 按 $\CK_{\gamma_{\rm BPS}}^{\Omega(\gamma_{\rm BPS}; u)}$ 跳变。我们还论证了：对充分大的 $R$，解确实存在。

\begin{yizhu}
原文的黎曼–希尔伯特问题由四要素确定：$\zeta \to 0, \infty$ 处的渐近行为（条件 f)，以及 e) 给出的共轭版本）、穿过 BPS 射线时的跳变（条件 g)、h)）、在 $\zeta$-平面上整体单值（由条件 a) 的乘性与上述跳变拼接而成），以及增长与全纯性要求（条件 b)、d)）。见原文 §2.2 条件 a)–f) 与式 (2.7)。
\end{yizhu}

这样的 $\CX^{\RH}_\gamma(\zeta)$ 与满足 a)–c)、d$'$)–f$'$) 的 $\CX_\gamma^{\vartheta}(\zeta)$ 之间有一个简单的对应。给定 $\CX^{\RH}_\gamma(\zeta)$，$\CX_\gamma^{\vartheta}(\zeta)$ 可由 $\CX^{\RH}_\gamma(\zeta)$ 在 $\zeta$ 中从射线 $\zeta \in e^{i \vartheta}\IR_+$ 出发作解析延拓得到。反之，给定 $\CX_\gamma^\vartheta(\zeta)$，我们把 $\zeta$-平面用 BPS 射线切成若干窄条，并规定 $\CX^{\RH}_\gamma(\zeta)$ 在包含射线 $\zeta \in e^{i \vartheta}\IR_+$ 的那个窄条上与 $\CX_\gamma^\vartheta(\zeta)$ 一致。

人们自然会猜想：若把这一手续施于本文所构造的 $\CX^\vartheta_\gamma(\zeta)$，所得的 $\CX^{\RH}_\gamma(\zeta)$ 将对一切 $R$ 给出 \cite{Gaiotto:2008cd} 中定义的黎曼–希尔伯特问题的一个解。然而，本文并不证明这一点。关键的问题在于：所得的 $\CX^{\RH}_\gamma(\zeta)$ 可能只是分片\emph{亚纯}的（在 $\zeta = 0, \infty$ 之外）。这与黎曼–希尔伯特问题不相容，因为其中要求 $\CX^{\RH}_\gamma(\zeta)$ 是分片\emph{全纯}的。在 $R$ 充分大时这不成问题：各个 $\CX^\vartheta_\gamma$ 的极点落在扇形区 $\IH_\vartheta$ 之外，不会出现在 $\CX^{\RH}_\gamma$ 中。重要的问题是：当我们走向小 $R$ 时，这些极点是否会移入围绕射线 $\zeta \in e^{i \vartheta}\IR_+$ 的那个窄条中。

尽管我们关于穿墙与 BPS 简并度的结果与这个问题无关，但它对用 \cite{Gaiotto:2008cd} 的方法描述\hk 度量一事是相关的。我们将在第 14 节% TODO 合稿改 \ref{sec:comparison}
中对此作稍为详细的讨论。
